This special session will bring together speakers from different disciplines under the umbrella of topological data analysis/vizualization. This will include mathematicians who study algorithms such as Mapper and concepts like persistent homology and dimension reduction, as well as researchers who are not mathematicians who have used these ideas successfully in their work. There will be ample discussion of software and opportunities to collaborate on larger TDA and visualization libraries, an area currently lacking in availability, user friendliness, and standardization. 

Mapper in particular has seen a variety of applications in various domains, such as neuroscience, clinical medicine, genomics, and sociology, each with their own approach to parameter selection and analysis. Despite this growing body of work, the field lacks a standardized methodology for parameter selection. Some work has been done on optimizing filter functions and covers, but their use for researchers in practice is necessarily limited, 
as the software landscape for Mapper is fragmented; implementations differ in their supported data formats, scalability, library APIs, and visualization capabilities. This leaves researchers unaware of best practices or critical technical details from other fields, which can lead to suboptimal, misleading, or even invalid analyses. The problem is not unique to Mapper; software for TDA and visualization is commonly out-of-date, unorganized, or closed-source.

The intended audience for this session is diverse; by bringing together mathematicians and researchers, we aim to bridge some of the gaps that exist between theory and practice, and foster new and exciting opportunities for further collaboration. The interdisciplinary nature of the session will bolster this, ideally attracting JMM attendees who might not otherwise have had the opportunity to see or work on such projects. This goes both ways; speakers who \textit{apply} TDA to their (non-mathematical) work can benefit from exposure of their ideas to the mathematical community, enabling them to gain partners and look at their work in new ways. Especially as the world becomes more and more saturated with data in the age of information and AI, there is more need than ever for mathematicians to make sense of complex, often high-dimensional data. TDA and visualization research enables both these mathematicians and those from other fields to make new connections and discoveries, and sessions such as these help keep those discoveries informed, accurate, and forward-focusing.